\chapter{Conclusion}

ARIA demonstrates a serious game approach to AI-assisted forensic training. The project does not present AI as either a perfect analyst or a useless tool. Instead, it places the player in a controlled investigation where AI output must be treated as a set of \textbf{hypotheses to test}.

The strongest design element is the combination of mixed true and hallucinated claims with an evidence review gate. This makes the desired professional behavior unavoidable: inspect artifacts first, then decide whether a claim belongs in the report. The fictional TechCorp fraud case gives the lesson a concrete shape through email, audio, video, PDF, and network evidence.

As a course artifact, the project delivers a working browser game, a deterministic scripted assessment mode, optional live AI experimentation, public deployment, documentation, and a debrief-oriented learning loop. Its next natural evolution would be scenario expansion and instructor-facing analytics, but the current prototype already provides a complete playable example of evidence-based AI calibration in digital forensics.
